\documentclass[12pt,a4paper]{article}

\usepackage[utf8]{inputenc}
\usepackage[T1]{fontenc}
\usepackage[spanish,es-tabla]{babel}
\usepackage{geometry}
\usepackage{booktabs}
\usepackage{longtable}
\usepackage{array}
\usepackage{hyperref}
\usepackage{enumitem}
\usepackage{graphicx}
\usepackage{xcolor}
\usepackage{float}
\usepackage{titlesec}

\geometry{margin=2.5cm}
\setlist[itemize]{noitemsep, topsep=3pt}
\setlist[enumerate]{noitemsep, topsep=3pt}
\titleformat{\section}{\large\bfseries}{\thesection.}{0.5em}{}
\titleformat{\subsection}{\normalsize\bfseries}{\thesubsection.}{0.5em}{}

\hypersetup{
    colorlinks=true,
    linkcolor=blue,
    urlcolor=blue,
    pdftitle={Proyecto Neo4j},
    pdfauthor={Grupo de trabajo}
}

\title{Proyecto No. 2 Neo4j\\\large Motor de recomendación de películas y series}
\author{Grupo de trabajo}
\date{Abril 2026}

\begin{document}

\maketitle
\tableofcontents
\newpage

\section{Resumen ejecutivo}

Este proyecto implementa un motor de recomendación de películas y series sobre Neo4j, utilizando AuraDB como motor de base de datos y Python con FastAPI como capa de carga y automatización. El enfoque principal es modelar un grafo con múltiples etiquetas, propiedades y relaciones, cargar más de 5000 nodos a partir de archivos CSV y exponer operaciones y consultas que permitan demostrar el funcionamiento del sistema.

\section{Objetivo}

Desarrollar una aplicación funcional que permita almacenar, consultar y manipular información de películas, usuarios e interacciones, garantizando la trazabilidad de los datos y el cumplimiento de la rúbrica del curso. El modelo de datos está orientado a facilitar recomendaciones basadas en similitud de usuarios, preferencias y contenido.

\section{Caso de uso}

El caso de uso seleccionado es el de motor de recomendación. La base de datos representa películas, géneros, directores, idiomas, colecciones y usuarios simulados con sus vistas, calificaciones, preferencias y relaciones de amistad. Con esta estructura se pueden ejecutar consultas de filtrado, agregación, recomendaciones y validación de conectividad del grafo.

\section{Tecnologías utilizadas}

\begin{itemize}
    \item Neo4j AuraDB.
    \item Python 3.13.
    \item Neo4j Python Driver.
    \item python-dotenv.
    \item FastAPI para la capa de backend.
    \item CSV como formato de carga de datos.
\end{itemize}

\section{Fuentes de datos}

Se partió de un archivo principal de películas y se generó data sintética adicional para usuarios e interacciones. Luego se realizó una limpieza ligera para normalizar campos, convertir tipos de datos y eliminar duplicados o valores inválidos.

\subsection{Archivos principales}

\begin{itemize}
    \item \texttt{data/movies.csv}
    \item \texttt{data/users.csv}
    \item \texttt{data/user\_views.csv}
    \item \texttt{data/user\_ratings.csv}
    \item \texttt{data/user\_preferences.csv}
    \item \texttt{data/user\_friendships.csv}
\end{itemize}

\subsection{Archivos limpios}

\begin{itemize}
    \item \texttt{data/clean/movies\_clean.csv}
    \item \texttt{data/clean/movie\_genres\_clean.csv}
    \item \texttt{data/clean/users\_clean.csv}
    \item \texttt{data/clean/user\_views\_clean.csv}
    \item \texttt{data/clean/user\_ratings\_clean.csv}
    \item \texttt{data/clean/user\_preferences\_clean.csv}
    \item \texttt{data/clean/user\_friendships\_clean.csv}
\end{itemize}

\section{Modelado del grafo}

\subsection{Etiquetas de nodos}

El modelo incluye las siguientes etiquetas:

\begin{itemize}
    \item Movie.
    \item User.
    \item Genre.
    \item Director.
    \item Language.
    \item Collection.
\end{itemize}

\subsection{Relaciones principales}

\begin{itemize}
    \item \texttt{DIRECTED}: director hacia película.
    \item \texttt{HAS\_GENRE}: película hacia género.
    \item \texttt{IN\_LANGUAGE}: película hacia idioma.
    \item \texttt{CONTAINS}: colección hacia película.
    \item \texttt{VIEWED}: usuario hacia película.
    \item \texttt{RATED}: usuario hacia película.
    \item \texttt{PREFERS}: usuario hacia género.
    \item \texttt{FRIEND\_OF}: relación entre usuarios.
\end{itemize}

\subsection{Propiedades clave}

\begin{itemize}
    \item Texto: títulos, nombres, comentarios, idiomas.
    \item Número: presupuesto, revenue, rating, progress, popularity.
    \item Booleano: premium.
    \item Listas: favorite\_languages, preferred\_genres.
    \item Fechas: release\_date, register\_date, view\_date, rating\_date.
\end{itemize}

\section{Limpieza y preparación de datos}

Se aplicó una limpieza ligera orientada a asegurar compatibilidad con Neo4j y consistencia durante la carga:

\begin{itemize}
    \item Normalización de espacios en blanco.
    \item Conversión de campos numéricos y fechas.
    \item Eliminación de duplicados exactos.
    \item Filtrado de identificadores inválidos.
    \item Control de rangos para puntuaciones y porcentajes.
\end{itemize}

\section{Carga de datos}

La carga se automatizó con Python usando el driver oficial de Neo4j. El proceso siguió este orden:

\begin{enumerate}
    \item Creación de constraints.
    \item Carga de nodos Movie, User, Director, Language y Collection.
    \item Carga de relaciones DIRECTED, HAS\_GENRE, IN\_LANGUAGE, CONTAINS, VIEWED, RATED, PREFERS y FRIEND\_OF.
    \item Validación posterior a la carga.
\end{enumerate}

\section{Resultados de la carga}

La carga final quedó con el siguiente resumen:

\begin{table}[H]
\centering
\begin{tabular}{lr}
\toprule
\textbf{Elemento} & \textbf{Cantidad} \\
\midrule
Movie & 4803 \\
User & 2200 \\
Genre & 19 \\
Director & 2350 \\
Language & 37 \\
Collection & 19 \\
\midrule
Total nodos & 9428 \\
Total relaciones & 107767 \\
\bottomrule
\end{tabular}
\caption{Conteo final de nodos y relaciones}
\end{table}

\begin{table}[H]
\centering
\begin{tabular}{lr}
\toprule
\textbf{Relación} & \textbf{Cantidad} \\
\midrule
DIRECTED & 4803 \\
HAS\_GENRE & 12113 \\
IN\_LANGUAGE & 4803 \\
CONTAINS & 12085 \\
VIEWED & 36304 \\
RATED & 19999 \\
PREFERS & 6668 \\
FRIEND\_OF & 21984 \\
\bottomrule
\end{tabular}
\caption{Conteo final de relaciones}
\end{table}

\section{Validación post carga}

La validación formal posterior a la carga se ejecutó exitosamente con estado PASS. El reporte generado confirmó lo siguiente:

\begin{itemize}
    \item No existen películas sin género.
    \item No existen películas sin director.
    \item No existen usuarios sin interacciones.
    \item No existen nodos aislados.
    \item El grafo es conexo desde el componente principal.
\end{itemize}

\subsection{Resumen de validación}

\begin{table}[H]
\centering
\begin{tabular}{lr}
\toprule
\textbf{Control} & \textbf{Resultado} \\
\midrule
movies\_without\_genre & 0 \\
movies\_without\_director & 0 \\
users\_without\_interactions & 0 \\
isolated\_nodes & 0 \\
\bottomrule
\end{tabular}
\caption{Controles de integridad}
\end{table}

\section{Estado actual del proyecto}

A la fecha de esta redacción, el proyecto ya cuenta con:

\begin{itemize}
    \item Data principal cargada y validada.
    \item Modelo extendido con labels \texttt{Language} y \texttt{Collection}.
    \item Data sintética para usuarios e interacciones.
    \item Script de limpieza de datos.
    \item Script de carga a Neo4j AuraDB.
    \item Script de validación post carga.
    \item API CRUD para nodos y relaciones.
\end{itemize}

\section{Pendientes inmediatos}

Los siguientes pasos serán:

\begin{enumerate}
    \item Definir y documentar consultas Cypher de demostración.
    \item Incorporar un algoritmo de recomendación (content-based o colaborativo) como evidencia de data science.
    \item Ampliar tipos de relaciones orientadas al caso de uso (por ejemplo \texttt{LIKED}, \texttt{WATCHLISTED}, \texttt{FOLLOWS\_DIRECTOR}, \texttt{CREATED}) para acercarse al objetivo del diseño.
    \item Preparar el guion de presentación final.
    \item Consolidar evidencias de pruebas de endpoints y consultas para la exposición.
\end{enumerate}

\section{Anexos}

\subsection{Rutas de scripts}

\begin{itemize}
    \item \texttt{scripts/generate\_synthetic\_users\_interactions.py}
    \item \texttt{scripts/preprocess\_data.py}
    \item \texttt{scripts/load\_clean\_data\_to\_neo4j.py}
    \item \texttt{scripts/validate\_post\_load.py}
    \item \texttt{scripts/fix\_post\_load\_orphans.py}
\end{itemize}

\subsection{Reporte de validación}

El detalle completo de la validación puede consultarse en \texttt{validation\_post\_load\_report.md}.

\end{document}
